\documentclass[ ../main.tex]{subfiles}
\providecommand{\mainx}{..}
\begin{document}
\section{Special cases and connections}
\label{sec:special}

The Bernoulli map model specializes to several important function classes.
In this section, we show how the model recovers the Bernoulli set, connects to binary channels in information theory, and accommodates constant functions and $n$-ary predicates.

\subsection{Indicator functions and Bernoulli sets}
\label{sec:indicator}

The indicator function of a set $\Set{A}$ over a universal set $\Set{U}$ is denoted by $\SetIndicator{\Set{A}} \colon \Set{U} \to \BitSet$.
A \emph{first-order} random approximate set of $\Set{A}$ with a false positive rate $\fprate$ and false negative rate $\fnrate$ is denoted by $\ASet{A}[\fprate][\fnrate]$, which may be modeled by a \emph{first-order} random approximate map $\APFun{\SetIndicator{\Set{A}}}[\fprate][\fnrate] \colon \Set{U} \to \BitSet$.

More generally, functions of type $\Set{U} \ATOverUnder{\mapsto}[\fprate][\text{\raisebox{3pt}{$\fnrate$}}] \BitSet$ model first-order random approximate subsets of the universal set $\Set{U}$.
This may also be indicated by $\AT{\RV{R}}[\fprate][\fnrate]$ given that $\RV{R}$ may realize any element in $\PS{\Set{U}}$.

\begin{remark}
The Bernoulli set model developed in the companion paper~\cite{aset} is thus a special case of the Bernoulli map model restricted to the codomain $\BitSet$.
All results for Bernoulli sets (error count distributions, confidence intervals, higher-order composition) apply to Bernoulli maps of indicator functions.
\end{remark}

\subsection{Bernoulli Booleans and binary channels}
\label{sec:bernoulli-bool}

We are not restricted to the set-indicator function. Consider the Bernoulli model applied to Boolean values.
How many functions are there of type $() \to \Bool$? There are two: \True and \False.
So, we can think of Boolean values as constant functions of type $() \to \Bool$.
When we apply the Bernoulli model to a Boolean value, we get a random approximate Boolean: the confusion matrix is a $2 \times 2$ matrix.
See \Cref{tab:confusion-bool}.

\begin{table}[ht]
\centering
\caption{Confusion matrix for the Bernoulli model applied to Boolean values.}
\label{tab:confusion-bool}
\begin{tabular}{lcc}
\toprule
latent/observed & \True & \False \\
\midrule
\True  & $p_1$   & $1-p_1$ \\
\False & $p_2$   & $1-p_2$ \\
\bottomrule
\end{tabular}
\end{table}

The confusion matrix in \Cref{tab:confusion-bool} has the maximum of
two degrees-of-freedom, since there are only two parameters.
For a \emph{pair} of independent Bernoulli Booleans---e.g., the joint output of two independent approximate maps---the observed outcome lies in $\Bool \times \Bool$, and the joint confusion matrix is $4 \times 4$.
A general $4 \times 4$ channel has $12$ free parameters, but element-wise independence forces the joint matrix to factor as a Kronecker product of two $2 \times 2$ matrices, reducing the count to $4$: one response probability per input class per Boolean component.
More generally, a map $f \colon X \to Y$ has its graph as a subset of $X \times Y$, so each element's membership test is binary and the Bernoulli independence axiom constrains the joint behaviour of $|Y|$ indicator variables per element.
Both cases can also be understood in information-theoretic terms.

\paragraph{Binary symmetric channel.}
In information theory~\cite{shannon1948,cover2006}, a binary channel transmits two
symbols, e.g., $0$ and $1$. The binary symmetric channel (BSC) flips bits with probability $\epsilon$. Its confusion matrix is
given in \Cref{tab:bsc}.

\begin{table}[ht]
\centering
\caption{Confusion matrix for the binary symmetric channel.}
\label{tab:bsc}
\begin{tabular}{lcc}
\toprule
latent/observed & \True & \False \\
\midrule
\True  & $1-\epsilon$  & $\epsilon$ \\
\False & $\epsilon$    & $1-\epsilon$ \\
\bottomrule
\end{tabular}
\end{table}

The BSC is the Bernoulli Boolean with $p_1 = 1 - \epsilon$ and $p_2 = \epsilon$.

\paragraph{Binary asymmetric channel.}
The binary asymmetric channel (BAC) flips 1's with a \emph{false negative}
rate of $\fnrate$ and 0's with a \emph{false positive} rate of $\fprate$.
See \Cref{tab:bac}.

\begin{table}[ht]
\centering
\caption{Confusion matrix for the binary asymmetric channel.}
\label{tab:bac}
\begin{tabular}{lcc}
\toprule
latent/observed & \True & \False \\
\midrule
\True  & $1-\fnrate$ & $\fnrate$ \\
\False & $\fprate$          & $1-\fprate$ \\
\bottomrule
\end{tabular}
\end{table}

The BAC is the Bernoulli Boolean with $p_1 = 1 - \fnrate$ and $p_2 = \fprate$.
The Bloom filter~\cite{bloom1970} set-indicator is the special case $\fnrate = 0$.

\subsection{Approximate constants}
\label{sec:constants}

The simplest non-trivial approximate map is the \emph{approximate constant}: a function $\Fun{f} \colon \Set{1} \to \Set{X}$ where $\Set{1}$ denotes a singleton set.
Then $\Fun{f}(1) \coloneqq y$ for some $y \in \Set{X}$, and the approximate constant $\APFun{f}(1)$ equals $y$ with probability $1-\fnrate$ and some other value in $\Set{X}$ with probability $\fnrate$.

\begin{remark}
Approximate constants establish that \emph{approximate values} are approximate maps from the unit type.
This observation is foundational for the algebraic data type treatment developed in the companion paper on random approximate algebraic types.
\end{remark}

\begin{theorem}[Kronecker factorization]
\label{thm:kronecker}
Let $\APFun{f} \colon \Set{X} \to \Set{Y}$ be a Bernoulli map with per-input independent errors and model order $n$ (so that $\Set{X}$ is partitioned into $n$ blocks within which all per-input channels are identical).
The joint confusion matrix $M$ over the function space $\Set{Y}^{\Set{X}}$ factors as
\begin{equation}
M = \bigotimes_{x \in \Set{X}} M_x,
\end{equation}
where each $M_x$ is the $|\Set{Y}| \times |\Set{Y}|$ per-input channel matrix at input $x$.
Since inputs within the same partition block share identical channel matrices, the number of free parameters is $n \cdot (|\Set{Y}| - 1)$, compared to $N(|\Set{Y}|^{|\Set{X}|} - 1)$ for a general $N$-th order channel on $\Set{Y}^{\Set{X}}$.
In the maximum-order case $n = |\Set{X}|$ (each input in its own block), this becomes $|\Set{X}| \cdot (|\Set{Y}| - 1)$.
\end{theorem}
\begin{proof}
By the element-wise independence assumption, the probability that $\APFun{f}$ maps the input vector $(x_1, \ldots, x_{|\Set{X}|})$ to the output vector $(y_1, \ldots, y_{|\Set{X}|})$ factors as
$\prod_{i} \Prob{\APFun{f}(x_i) = y_i \mid f(x_i)}$.
Each factor is an entry of the per-input channel matrix $M_{x_i}$.
The joint matrix is therefore the Kronecker product of the per-input matrices.
Since the i.i.d.\ constraint within each of the $n$ partition blocks forces all inputs in a block to share the same channel matrix, the free parameters are determined by $n$ distinct channel matrices, each with $|\Set{Y}| - 1$ free parameters (one row of a stochastic matrix), giving $n(|\Set{Y}| - 1)$ total.
\end{proof}

\begin{remark}[Function-space viewpoint]
\label{rem:function_space}
\Cref{thm:kronecker} shows that every Bernoulli approximation is an instance of the same phenomenon: a latent function $\Fun{f'} \colon \Set{X} \to \Set{Y}$ is approximated by an observed function $\Fun{f} \colon \Set{X} \to \Set{Y}$, where the per-input errors are independent and the joint confusion matrix factors as a Kronecker product.

This gives a unified account of the approximation hierarchy:
\begin{itemize}
\item \textbf{Bernoulli Boolean}: $\Set{X} = \Set{1}$, $\Set{Y} = \Bool$.
    The function space has $2^1 = 2$ elements (\True and \False), and the confusion matrix is $2 \times 2$ with at most $2$ partition blocks and $2$ free parameters.
\item \textbf{Bernoulli set}: $\Set{X} = \Set{U}$, $\Set{Y} = \BitSet$.
    The indicator function $\SetIndicator{\Set{A}} \colon \Set{U} \to \BitSet$ is approximated per-element, and each element's $2 \times 2$ channel contributes one parameter per partition block.
\item \textbf{Bernoulli map}: general $\Set{X}$, $\Set{Y}$.
    For each input $x \in \Set{X}$, the per-input channel is $|\Set{Y}| \times |\Set{Y}|$ with at most $|\Set{Y}| - 1$ free parameters per partition block.
    The full confusion matrix factors as a Kronecker product of these per-input channels.
\item \textbf{Bernoulli relation}: $\Set{X} = \Set{A} \times \Set{B}$, $\Set{Y} = \BitSet$.
    A relation is a set of pairs, so it is a Bernoulli set on the product domain $\Set{A} \times \Set{B}$.
\end{itemize}
In every case, the Bernoulli independence axiom reduces the parameter count from exponential (in $|\Set{X}|$) to linear.
The Boolean is the atomic case; all others are Kronecker products of Boolean-like channels across the input domain.
\end{remark}

\subsection{Relations as $n$-ary predicates}
\label{sec:nary}

A $k$-ary predicate function $\Fun{f} \colon \Set{T}_1 \times \cdots \times \Set{T}_k \to \Bool$ maps $k$-tuples to \True or \False.
A $k$-ary predicate is equivalently a $k$-ary \emph{relation}: it determines which tuples are ``related.''
For instance, an ordered set answers questions like $a \leq b$ by checking whether $(a,b)$ belongs to the binary relation ${\leq}$.

When such predicates are approximated, we may chain the tests.
For example, given two approximate binary predicates $\APFun{f}[\fprate_1][\fnrate_1] \colon \Set{X} \times \Set{Y} \to \BitSet$ and $\APFun{g}[\fprate_2][\fnrate_2] \colon \Set{Y} \times \Set{Z} \to \BitSet$, a $3$-ary predicate is approximated by
\begin{equation}
	\Fun{w}(x,y,z) \coloneqq \Fun{f}(x,y) \land \Fun{g}(y,z)\,.
\end{equation}

\begin{remark}
The full treatment of approximate $n$-ary relations, including join, project, and select operations, is developed in the companion paper on approximate relations.
\end{remark}

\end{document}
